\documentclass{article}
\usepackage{iclr2025_conference,times}
\input{math_commands.tex}

\usepackage{amsmath,amssymb, comment}
\usepackage{graphicx}
\usepackage{float}
\usepackage{xcolor}
\usepackage{url}
\bibliographystyle{unsrtnat}

\iclrfinalcopy          % non-anonymous (shows authors); comment out for anonymous submission

% gray-italic scaffolding note: what should go in a section
\newcommand{\scaffold}[1]{{\par\color{gray}\itshape #1\par\medskip}}

% placeholder box for figures still to be generated
\newcommand{\figplaceholder}[2][0.6\linewidth]{%
  \fcolorbox{gray}{gray!8}{\parbox[c][3.2cm][c]{#1}{\centering\ttfamily\footnotesize
  \textcolor{red}{[GENERATE]}\\[2pt] #2}}}

\title{Tools as Resources: Measuring Online Tool-Investment Capability in LLM Agents}

\author{%
Roshan Klein-Seetharaman\thanks{Equal contribution (co-first authors).} \quad
Daniel Wang\footnotemark[1] \quad
Andrew Xu\footnotemark[1] \\
Sea12 Technologies \quad Yale University \\
\texttt{daniel.wang.dzw22@yale.edu}}

\newcommand{\fix}{\marginpar{FIX}}
\newcommand{\new}{\marginpar{NEW}}

\begin{document}
\maketitle
\begin{abstract}
Creating a reusable tool is an investment: an agent pays a fixed cost now in
exchange for uncertain future reuse. We introduce a paired benchmark that tests
whether language-model agents preserve online allocation behavior when the same
investment decision moves from an abstract keep/pass frame to reusable code
construction. Claude Haiku reserves budget in the abstract frame but makes
100\% of its script builds on first sight, compared with 28\% abstract
first-sight commitment (paired difference $+72.2$ percentage points, 95\%
seed-bootstrap interval $[+61.1,+83.3]$). A preregistered framing ladder
localizes this shift to required code emission: a declarative solver claim
yields 39\% first-sight commitment, while requiring code---without
correctness-dependent reward or a hand-solving option---raises it to 100\%
(paired difference $+61.1$ points, $[+38.9,+83.3]$). Across build charges
spanning eager, selective, and never-build regimes, Haiku's abstract choices
adjust to the economics, whereas its code-required choices remain at
98--100\% first sight, including when never building is optimal. The shift also appears in Claude
Opus and, after a preregistered competence gate, in GPT-5.4-mini (16.7\%
abstract versus 87.9\% code-required). In contrast, competence-qualified
GPT-5.6 Sol preserves reserve across the same comparison (19.4\% versus
27.8\%; the paired interval includes zero), demonstrating that the tested
capability is achievable rather than structurally precluded by code generation.
In a separate proof-of-possibility case study, a Qwen-14B policy trained from
reward acquires abstract reserve and generalizes it across held-out lexical
reskins, but remains eager in reusable script creation.
Together, these results establish online tool investment as a discriminative
capability boundary: constructive framing can suppress an available allocation
policy, but a capable model can preserve it.
\end{abstract}

\section{Introduction}

Language-model agents increasingly create scripts, wrappers, macros, and skills
that can be reused across tasks. Existing evaluations ask whether a model can
write a correct tool, choose an available tool, or benefit from a growing
library \cite{cai2024latm,qian2023creator,huang2024ultratool,qu2025toollearningsurvey}.
They largely leave a prior decision unmeasured: \emph{when should the agent
build?}

Building is worthwhile only when its fixed cost can be amortized over future
reuse. An agent facing an unfamiliar problem type therefore has three
possibilities: build immediately, wait for evidence of recurrence, or never
build. Building every initially plausible tool wastes a scarce construction
budget on one-offs; waiting too long forfeits useful reuse. This is an online
rent-or-buy problem, related to ski rental \cite{karlin1994competitive}, but
with multiple latent problem classes, a shared build budget, and a constructive
action generated by a language model.

Eager building is ambiguous when observed in an ordinary tool benchmark. The
agent may lack an allocation policy, may possess it but fail to activate it in
the tool context, or may understand the economics while being pulled toward
the current problem. We separate these explanations by presenting identical
problem streams and information in paired decision frames. The abstract frame
asks the model to keep or pass colored items. The constructive frame asks it
to claim or write a reusable solver for the corresponding numeric problem
class. This design first establishes whether an allocation policy is present,
then tests whether it survives constructive reframing.

Our results support five conclusions. First, with the number of latent classes
disclosed in both frames, Haiku commits on first sight in 28\% of abstract
decisions but 100\% of reusable-script decisions. Second, a preregistered
framing ladder and a construction-only control isolate required code emission
as sufficient to reproduce the full shift within this benchmark. Third, the
affected code-required policy is economically inelastic: it remains almost
fully eager as visible build charges move from zero to a level at which every
build is net-negative. Fourth, the benchmark is difficult but passable. Opus and
GPT-5.4-mini exhibit the framing shift, whereas GPT-5.6 Sol preserves reserve
after code becomes required. Fifth, as a convergent case study, a Qwen-14B
policy can acquire reserve from reward and preserve it when balls and urns are
replaced by unseen abstract vocabularies, yet remains as eager as its base model
during script construction.

The paper contributes:
\begin{enumerate}
    \item an online tool-investment formulation and benchmark with uncertain
    recurrence, scarce irreversible builds, paired streams, and explicit
    timing outcomes;
    \item a same-information test that separates absent allocation competence
    from framing-dependent non-expression;
    \item a causal study combining a preregistered framing ladder, a
    code-required control, and an economic response surface;
    \item a preregistered, competence-gated model panel that contains both a
    large cross-frame shift and successful cross-frame preservation; and
    \item a proof-of-possibility reward-learning case study distinguishing
    policy acquisition from cross-context activation.
\end{enumerate}

\section{Related Work}

\paragraph{Tool creation and skill libraries.}
LATM amortizes an expensive model's tool creation over many uses
\cite{cai2024latm}; CREATOR separates tool creation, decision, execution, and
rectification \cite{qian2023creator}; and ReGAL, LILO, and TroVE induce reusable
program abstractions or libraries \cite{stengeleskin2024regal,grand2024lilo,wang2024trove}.
These systems generally assume that a recurring task warrants construction or
induce libraries from a pre-existing corpus. Work on library-learning failures
shows that nominally reusable abstractions may see little actual reuse
\cite{berlotattwell2024librarylearning}. Library Drift studies the downstream
cost of unbounded accumulation and repairs it through retirement and bounded
library management \cite{zhang2026librarydrift}. Our object precedes these:
whether to author an irreversible artifact before recurrence is known.

\paragraph{Tool-use evaluation.}
UltraTool and broader surveys evaluate planning, tool creation, selection, and
execution \cite{huang2024ultratool,qu2025toollearningsurvey}. These are primarily
capability or accuracy evaluations. Our outcome is a disposition conditional
on demonstrated competence: whether an agent preserves budget until reuse
evidence justifies construction.

\paragraph{Sequential decision-making and early commitment.}
Prior work evaluates regret of LLMs selecting from fixed actions in online
learning and games \cite{park2025regret}. Meta-bandit training can produce
effective but increasingly exploitative policies that abandon exploration too
early \cite{chen2025greedy}. Mehta \cite{mehta2026commit} studies
representational commitment within a single agent trajectory. Our commitment
is instead an externally visible, irreversible construction decision whose
cost is amortized over a stream. We borrow the rent-or-buy perspective from
online algorithms \cite{karlin1994competitive}, but do not claim that our
multi-class, budgeted benchmark is identical to textbook ski rental.

\section{Online Tool Investment}

\subsection{Decision process}

An episode contains a horizon of $T$ items with latent classes
$c_t \in \{1,\ldots,N\}$. At time $t$, the agent observes the current item and
history $h_t$, has remaining build budget $b_t$, and chooses whether to commit
to class $c_t$. A commitment is irreversible and consumes one of $B$ available
builds. Once a class has been built, later occurrences can reuse its artifact.
The class frequencies are unknown at the start of the episode.

For the economic experiments, a solved or collected occurrence contributes
one point and a commitment costs $K$. For policy $\pi$,
\begin{equation}
U(\pi) =
\sum_{t=1}^{T} r_t(\pi)
- K\sum_{c=1}^{N}\mathbb{I}\{\pi\text{ commits to }c\},
\end{equation}
subject to at most $B$ commitments. The best action depends on both recurrence
evidence and remaining budget. A one-off class should not consume a build;
a sufficiently frequent class should.

\subsection{Benchmark construction}

We instantiate the process with streams of $T=60$ numeric problems drawn from
eight latent problem families. Each canonical stream contains approximately
three ``hot'' classes that recur 13--23 times and several ``trap'' classes that
appear only 1--6 times. An early-trap constraint places a rare class among the
first $B$ opportunities, making first-sight commitment costly. The hardened
numeric pool is designed to make unaided hand solving approximately unavailable
while retaining script-based solvability.

The abstract and constructive environments are views of the same latent class
sequence. In the abstract urn frame, the model observes a colored item and
chooses \textsc{keep} or \textsc{pass}. Keeping a class commits budget and
collects its occurrences. In the reusable-script frame, the model encounters a
numeric problem and may write a class-level solver that can be run on later
instances. We pair identical latent class identities, roles, counts, and arrival
orders across conditions. Numeric payloads may be rendered at a model-calibrated
magnitude without changing that latent sequence.

The main causal condition, A2, tells both the model and the reference that
exactly $N=8$ classes occur. A separate no-$N$ condition is ecologically
natural but gives the Bayesian reference more information than the model.
Unless stated otherwise, causal comparisons use A2.

\subsection{Outcomes and references}

Our primary outcomes are behavioral timing measures:
\begin{itemize}
    \item \textbf{First-sight commitment}: the proportion of realized
    commitments made on a class's first occurrence.
    \item \textbf{First-sight hazard}: commitments on eligible first-occurrence
    turns divided by eligible first-occurrence turns.
    \item \textbf{Lateness}: the occurrence index at commitment minus the
    class's first occurrence; zero denotes immediate commitment.
\end{itemize}

For the original urn experiment, we use an exact finite-horizon Bayesian policy
with a symmetric Dirichlet prior as a labeled behavioral comparator. It is
Bayes-optimal for that prior, not for the engineered generator. For economic
regret we instead use an exact hindsight net optimum. On a realized stream, a
class of size $s_c$ built at first sight contributes $s_c-K$; the hindsight
reference chooses up to $B$ classes with the largest positive values. This
reference is prior-free and provides an offline upper bound, not an online
policy available to the model.

\section{Experimental Protocol}

The primary controlled subject is Claude Haiku 4.5. Claude Opus 4.8 provides a
frontier-model paired test. Qwen2.5-Coder-14B is used for reward learning.
GPT-5.4-mini and GPT-5.6 Sol are evaluated under a preregistered competence gate
with matched R0 and R2c conditions and \texttt{reasoning\_effort=none}.

Core Haiku comparisons use 12 canonical streams (seeds 2000--2011).
Repeated conditions on these streams are not independent samples: the effective
stream sample is 12, not the number of sessions. The final Qwen transfer test
uses 24 paired held-out streams (2000--2023), disjoint from training seeds.
The Opus script arm re-renders the same 12 canonical latent streams at magnitude
1000, which makes the tunable numeric families hand-hard for Opus while
preserving every class identity, role, count, and arrival position. A preflight structural
projection assertion passes on all 12 seeds. Josephus remains partly
hand-solvable for Opus at any magnitude; we therefore treat Opus timing as the
primary paired outcome, report a Josephus-excluded sensitivity, and do not use
its script-arm regret as a payoff-isomorphic cross-frame estimate.
We report numerators and denominators, pooled arm summaries, and mean paired
seed differences. Unless stated otherwise, 95\% intervals resample paired
seeds with replacement 10,000 times and take the 2.5th and 97.5th percentiles
of the bootstrap distribution. For first sight among realized commitments, a
seed with zero commitments has an undefined proportion and is excluded from
that estimand; first-sight hazard, commitment count, and zero-commit incidence
remain defined and are reported alongside it. Economic interactions are
computed within seed before resampling.

Mechanical validity is checked independently of behavioral outcomes. The
harness asserts stream hashes, separates transport failures from model
decisions, records malformed and unknown tool calls, and checks token-cap
truncation. Only mechanical failures permit reruns. Prompt, tool-choice, and
token-cap repairs discovered during benchmark development are documented in
the artifact record.

\section{Same Information, Different Allocation}

\subsection{Haiku's paired dissociation}

Figure~\ref{fig:task-core} illustrates the paired task construction and the
main A2 contrast. Table~\ref{tab:haiku2x2} presents the same 12 streams under both information
conditions. Disclosing $N$ makes Haiku substantially more selective in the
abstract frame: first-sight commitment falls from 53\% to 28\%, and lateness
rises from 0.86 to 1.19. The same disclosure has no effect in the script frame:
all builds occur on first sight under both conditions.

\begin{figure}[t]
\centering
\includegraphics[width=\linewidth]{figures/fig1_task_and_core_dissociation.png}
\caption{\textbf{One allocation problem, two decision frames.}
\textbf{(a)} The abstract and script environments expose the same latent class
stream; colors and numeric problem families are alternative surface
representations. \textbf{(b)} On 12 paired A2 streams, Haiku reserves in the
abstract R0 frame but constructs scripts at first sight in R3. Thin lines show
paired stream-level values and large points show pooled first-sight
commitment. Both conditions disclose exactly $N=8$ latent classes.}
\label{fig:task-core}
\end{figure}

\begin{table}[t]
\centering
\caption{Haiku allocation under paired information and framing conditions.
Timing is the primary outcome; the Bayesian comparator is prior-dependent.}
\label{tab:haiku2x2}
\begin{tabular}{lrrrr}
\hline
 & Urn no-$N$ & Urn A2 & Tool no-$N$ & Tool A2 \\
\hline
First-sight commitment (\%) & 53 & 28 & 100 & 100 \\
Mean lateness & 0.86 & 1.19 & 0.00 & 0.00 \\
Regret vs.\ Bayes comparator & $640$ & $0$ & $1308$ & $1633$ \\
Standard error & $285$ & $344$ & $520$ & $553$ \\
\hline
\end{tabular}
\end{table}

The A2 comparison is the key identification result. Haiku receives the same
class-count information and acts on the same class streams, yet waits in the
abstract task and commits immediately when the action is reusable script
construction. The mean paired R3-minus-R0 first-sight difference is $+72.2$
percentage points (95\% seed-bootstrap interval $[+61.1,+83.3]$). This is
evidence for framing-dependent non-expression, not a claim that all cognition
is held constant.

\subsection{Qualitative behavior}

Abstract decisions commonly refer prospectively to recurrence and preserving
budget. Tool-frame rationales instead center on the current problem: the model
writes a script to solve the instance in front of it, notices budget only after
spending it, and reuses scripts reactively if a class later returns. These
transcripts are consistent with the timing result but are not themselves the
causal evidence.

\section{What Triggers Eager Construction?}

\subsection{Preregistered framing ladder}

We vary the decision representation in four paired rungs:
\begin{description}
    \item[R0:] abstract free-text \textsc{keep}/\textsc{pass};
    \item[R1:] the same abstract decision through tool calls;
    \item[R2:] real numeric content with a declarative solver claim, but no code;
    \item[R3:] full reusable-script construction.
\end{description}
The preregistered localization criterion was the earliest adjacent increase of
at least 30 percentage points in first-sight commitment. The ladder is coarse,
not a complete factorial design.

\begin{table}[t]
\centering
\caption{Haiku framing ladder on 12 paired A2 streams.}
\label{tab:ladder}
\begin{tabular}{lrrrrr}
\hline
 & R0 & R1 & R2 & R2c & R3 \\
\hline
First sight (\%) & 28 & 17 & 39 & 100 & 100 \\
First-sight count & 10/36 & 6/36 & 14/36 & 36/36 & 34/34 \\
Mean lateness & 1.194 & 2.167 & 0.833 & 0.000 & 0.000 \\
\hline
\end{tabular}
\end{table}

R1 changes action modality without reproducing eagerness. R2 introduces real
problem content and a current-plus-future commitment, increasing first-sight
behavior by 22 points relative to R1 but remaining below the locked threshold.
The mean paired adjacent differences (95\% seed-bootstrap intervals) are
R1--R0 $=-11.1$ points $[-30.6,+8.3]$, R2--R1 $=+22.2$
$[+2.8,+44.4]$, R2c--R2 $=+61.1$ $[+38.9,+83.3]$, and R3--R2c
$=0.0$ $[0.0,0.0]$. These descriptive intervals do not replace the
preregistered 30-point localization rule and are not multiplicity-adjusted; the
final degenerate interval reflects ceiling observations in both rungs, not a
general equivalence claim.
Figure~\ref{fig:ladder} displays the paired stream trajectories and the
construction-only R2c control.

\begin{figure}[t]
\centering
\includegraphics[width=\linewidth]{figures/fig2_framing_ladder.png}
\caption{\textbf{Preregistered framing ladder and construction-only control.}
Thin lines connect the same 12 canonical streams across conditions; large
points show pooled first-sight commitment. Tool-call modality (R1) and a
declarative solver claim (R2) do not reproduce full eagerness. Requiring code
while preserving R2's payoff and removing correctness stakes (R2c) raises
first-sight commitment from 39\% to 100\%, matching full script construction
(R3).}
\label{fig:ladder}
\end{figure}

\subsection{Construction-only control}

R2c retains R2's binary claim/skip loop, budget, streams, payoff, and lack of a
hand-solving option. Its only addition is a required non-empty
\texttt{solve(inputs)} implementation. Payoff is deliberately independent of
code correctness, so correctness incentives cannot explain a timing shift.
R2c reaches 100\% first-sight commitment and zero lateness, exactly matching
R3 (Table~\ref{tab:ladder}). Within this benchmark, required code emission is
therefore sufficient to reproduce the observed effect; tool-call modality,
correctness stakes, and a hand-solving escape valve are not necessary.

Code quality provides a secondary diagnostic. Only 6 of 392 tested R2c
instances (2\%) are correct, and 31 of 36 claimed solvers are degenerate across
their class. Inspection shows structurally plausible code that hard-codes
incidental parameters from the current instance instead of using the input
fields. This suggests current-instance fixation, but correctness is post hoc
and does not enter the causal R2/R2c comparison.

\section{Economic Sensitivity}

We cross abstract R0 and code-required R2c with budgets
$B\in\{1,3,5\}$ and visible build charges $K$. The main plotted charges are
$K=0$, where eager construction of hot classes can pay; $K=10$, where hot
classes pay but traps do not; and $K=24$, where the largest class contains only
23 occurrences and every build is therefore net-negative. The $K=0,20,24$
panel contains 216 sessions, and the added $K=10$ selective panel contains 72.
All sessions pass stream, parsing, and token-cap checks.

\begin{table}[t]
\centering
\caption{First-sight commitment hazard (\%) over the economic surface. Entries
within a cell correspond to $B=1/3/5$.}
\label{tab:econ}
\begin{tabular}{lrrr}
\hline
Policy/frame & $K=0$ & $K=10$ & $K=24$ \\
\hline
R0 & 17/25/42 & 0/6/16 & 0/9/5 \\
R2c & 100/95/100 & 100/100/100 & 100/97/98 \\
Hindsight optimum & 33/61/83 & 33/57/39 & 0/0/0 \\
\hline
\end{tabular}
\end{table}

Figure~\ref{fig:economic} shows the interaction directly. R0 responds to price.
At $K=10$, it waits 1.56--3.75 occurrences on average,
and its first-sight hazard is at most 16\%. At $K=24$, it remains imperfect but
is far less eager than at zero charge. R2c, by contrast, stays at 98--100\%
first-sight across the surface. It commits immediately in the selective regime
and spends its budget even when every build is provably dominated by doing
nothing. At $B=5,K=24$, regret against the exact hindsight optimum is
$69.9\pm2.1$ for R2c and $60.3\pm6.0$ for R0.

The interaction matters more than either policy's absolute optimality. R0 is
too cautious at $K=0$ and still somewhat too eager at high charges, but its
timing changes with the visible economics. R2c's timing barely changes. For
Haiku, the code-required frame largely removes behavioral elasticity to the
investment price. We quantify this with a seed-paired endpoint interaction:
the $K=0$-to-$24$ reduction in first-sight hazard under R0 minus the
corresponding reduction under R2c. The mean interaction is $+43.3$ points
$[+18.3,+68.4]$ for $B=1$, $+21.4$ $[+7.9,+36.1]$ for $B=3$, and
$+37.0$ $[+19.8,+55.2]$ for $B=5$ (95\% seed-bootstrap intervals).
Because R2c is near ceiling throughout, this interaction is driven mainly by
R0's price response. Appendix~\ref{app:economic} reports hazard, commitments,
zero-commit incidence, lateness, net utility, and regret together.

\begin{figure}[t]
\centering
\includegraphics[width=\linewidth]{figures/fig3_economic_elasticity.png}
\caption{\textbf{Abstract allocation responds to price; code-required
commitment does not.} First-sight commitment hazard is shown across visible
build charges for each build budget. The abstract R0 policy becomes more
selective as construction grows expensive, broadly co-moving with the exact
hindsight net optimum. R2c remains near 100\% first sight, including at
$K=24$, where every build is net-negative. Each cell reuses the same 12 paired
Haiku streams; the hindsight reference is an offline upper bound, not an
online policy.}
\label{fig:economic}
\end{figure}

\section{A Difficult but Achievable Capability}

The paired benchmark admits three distinguishable states: policy absence,
policy possession with constructive suppression, and cross-frame preservation.
The model panel is summarized in Table~\ref{tab:models} and
Figure~\ref{fig:capability-map}.

\begin{table}[t]
\centering
\caption{Competence-gated framing comparisons. Arm rates are pooled ratios;
differences are means of paired seed-level differences. Opus uses the matched
A2 urn/tool endpoints.}
\label{tab:models}
\begin{tabular}{lrrrr}
\hline
Model & Abstract first sight & Code/tool first sight & Difference & Interpretation \\
\hline
Haiku 4.5 & 28\% & 100\% & +72.2 pp & Suppressed \\
Opus 4.8 & 0\% & 100\% & +100 pp & Suppressed \\
GPT-5.4-mini & 16.7\% & 87.9\% & +75.8 pp & Suppressed \\
GPT-5.6 Sol & 19.4\% & 27.8\% & +8.3 pp & Preserved \\
\hline
\end{tabular}
\end{table}

\begin{figure}[t]
\centering
\includegraphics[width=0.78\linewidth]{figures/fig4_capability_map.png}
\caption{\textbf{Online tool investment is difficult but achievable.} The
diagonal denotes preservation of first-sight timing across abstract and
code-required frames. Haiku, Opus, and GPT-5.4-mini lie far above it: each
reserves abstractly but becomes eager during construction. GPT-5.6 Sol remains
near the diagonal, demonstrating successful cross-frame preservation. The
hollow Qwen-14B point is descriptive only: its urn and final script values come
from different matched evaluation panels and should not be compared as a
competence-gated model estimate.}
\label{fig:capability-map}
\end{figure}

Opus reserves completely in the urn but uses all of its script budget at first
sight; its paired difference is $+100$ points, with a degenerate seed-bootstrap
interval $[+100,+100]$ because every observed seed has the same endpoint
values. All 36 script builds occur at first sight with zero lateness. Josephus
accounts for three of those builds; excluding it leaves 33/33 first-sight builds
and zero lateness. Thus the result is not driven by the one family for which
hand solving remains a partial option. The degenerate interval is not a
population-level certainty claim. GPT-5.4-mini passes
the preregistered abstract gate, then rises from
16.7\% R0 to 87.9\% R2c first-sight commitment; the paired difference is
$+75.8$ points with a 95\% bootstrap interval of $[+54.5,+93.9]$. Seed-level
heterogeneity remains: one R2c seed waits substantially and one makes no
commitments. The displayed arm rates are pooled ratios: $29/33-6/36=71.2$
points. The $+75.8$-point inferential estimand instead averages the 11
within-seed differences for which both among-commitment proportions are
defined; the zero-commit R2c seed is excluded from that conditional estimand.
R0 and R2c hazards are 8.7\% and 59.2\%, respectively; this unconditional
analysis retains the zero-commit seed.

GPT-5.6 Sol also passes the abstract gate, collecting 98.5\% of the Bayesian
comparator's payoff. Unlike the affected models, it remains selective when code
is required: R0 and R2c first-sight rates are 19.4\% and 27.8\%, with mean
lateness 1.14 and 1.11. The paired difference is $+8.3$ points with a 95\%
interval of $[-11.1,+27.8]$. Under the locked protocol, Sol preserves the
timing policy across construction.

This result is positive evidence for the benchmark. Code generation does not
mechanically force eager commitment, and the task is not an impossible trap.
It can distinguish a model that integrates allocation with construction from
models that demonstrate allocation only abstractly. The two GPT points do not,
however, establish that scale or generic capability causes preservation.
Post-training, architecture, reasoning policy, and tool-use experience remain
possible explanations.

\section{Convergent Case Study: Acquisition Without Constructive Activation}

As a separate proof-of-possibility case study, we ask whether reward can install
reserve in a model that is initially eager in both frames, and whether the
learned policy then activates during script creation. Starting from untouched
Qwen2.5-Coder-14B-Instruct, one QLoRA \cite{dettmers2023qlora} per-decision PPO
\cite{schulman2017ppo} run with no demonstrations reduces first-sight
commitment on 24 held-out abstract streams from 75\% to 32\%; lateness rises
from 0.375 to 0.903, and collected items rise from 87\% to 101\% of the
Bayesian comparator. The critic receives true class-rate features unavailable
to the actor \cite{pinto2018asymmetric}, and only one successful run is
analyzed. The result therefore establishes acquisition in this configuration,
not robust or sample-efficient learning. Appendix~\ref{app:rl} gives the
training procedure and pilot history.

The learned reserve is not tied to the literal words \emph{ball},
\emph{color}, and \emph{bag}. On three unseen lexical reskins that preserve the
abstract keep/pass structure, RL-final makes 19\% of commitments on first sight,
versus 89\% for the base model, with mean lateness 0.968 versus 0.134. This is
lexical generalization, not general task transfer.

Transfer nevertheless remains context-bound. On 24 matched q8\_0 reusable-script
streams, RL-final and base make 95\% of builds on first sight (58/61 and 59/62).
Expressing the abstract reskins through zero-argument keep/pass tool calls
produces partial, vocabulary-sensitive transfer (62\% for RL-final versus 99\%
for base). Thus the learned disposition survives new words but degrades under a
modality change and is absent in reusable construction
(Figure~\ref{fig:rl-transfer}). Detailed modality probes, per-vocabulary
results, and retry diagnostics are deferred to Appendix~\ref{app:rl}.

\begin{figure}[t]
\centering
\includegraphics[width=0.82\linewidth]{figures/fig5_rl_transfer_boundary.png}
\caption{\textbf{Reward-learned reserve remains context-bound.} The trained
Qwen-14B policy reserves in the original abstract urn and across three lexical
reskins. Transfer weakens when the same keep/pass decision is expressed through
tool calls and disappears in reusable script construction, where base and
RL-final are both 95\% first sight. Values are pooled publication results; the
script point is the matched q8\_0, $n=24$, EFR4 rerun. The four columns are
distinct evaluation panels rather than one paired trajectory.}
\label{fig:rl-transfer}
\end{figure}

\section{Discussion}

\paragraph{Possession, acquisition, and activation.}
The experiments separate three questions that ordinary tool benchmarks
conflate. A policy may be absent, as in baseline Qwen; acquired from reward but
context-bound, as in the Qwen case study; or already present but suppressed by
constructive framing, as in Haiku, Opus, and GPT-5.4-mini. GPT-5.6 Sol shows
that activation during construction is itself a measurable capability.

\paragraph{Why might construction suppress allocation?}
Required code emission may invoke a reactive coding-assistant policy centered
on solving the visible instance. It may also collapse ``solve now'' and
``invest for later'' into one salient action. R2c rules out correctness stakes
and hand-solving availability as necessary causes, but it does not identify a
neural mechanism. The transcript and code-quality evidence supports
current-instance fixation only as a bounded behavioral interpretation.

\paragraph{Implications for agents.}
Tool-building systems should separate the decision to invest from execution of
the build. An explicit recurrence ledger, remaining-budget state, or evidence
threshold could make the investment decision legible before code generation
begins. These are design implications, not tested interventions. Evaluation
should score delayed utility and library opportunity cost rather than only
immediate task success or tool correctness.

\paragraph{Implications for capability evaluation.}
A useful online-investment evaluation needs two conditions. An abstract
competence gate establishes that the model can reserve at all; the constructive
condition tests whether timing survives realistic action. Either alone is
ambiguous. Future work should test whether cross-frame preservation predicts
the quality of libraries and long-horizon resource allocation in naturalistic
repository, data-cleaning, API, and analysis workloads.

\section{Limitations}

The benchmark uses one synthetic numeric environment with engineered hot and
trap recurrence. Real tasks have ambiguous class boundaries, nonstationary
frequencies, partial reuse, and variable construction quality. R2c intentionally
decouples payoff from correctness to isolate code emission; it is not a
realistic model of artifact value.

The core causal ladder is centered on Haiku and uses 12 independent canonical
streams. Hundreds of crossed sessions do not increase the independent stream
sample beyond 12. The model panel establishes heterogeneity and achievability,
not a scaling trend or broad population estimate. The full economic response
surface is run only on Haiku, so GPT-5.6 Sol's result demonstrates cross-frame
timing preservation, not economic optimality at varying prices.

The Bayesian urn comparator is prior-dependent and mismatched to the engineered
generator. We therefore treat timing as primary and use the exact hindsight
net optimum only for realized economic regret. That optimum is unavailable to
an online model and should not be mistaken for an implementable policy.

The RL evidence contains one successful run and a privileged critic. It does
not show robust retraining, non-privileged learnability, or sample efficiency.
Finally, prompt and tool-choice formats affect deliberation. The framing ladder
is coarse rather than fully factorial, and the R1-to-R2 contrast bundles action
labels with problem content. Harness repairs were logged transparently, but
occurred during benchmark development.

\section{Conclusion}

Reusable tool creation is not only a coding problem; it is an online investment
decision under uncertain reuse. Paired experiments show that several models
which reserve abstractly commit immediately when required to emit code. In
Haiku, required code emission reproduces the full shift and renders commitment
nearly invariant to visible economics. A learned Qwen reserve policy likewise
remains context-bound. Yet GPT-5.6 Sol preserves reserve across the locked
framing change, demonstrating that online allocation during constructive action
is difficult but achievable. Online tool investment is therefore a distinct
capability worth measuring, not a universal model failure.

\bibliography{references}

\appendix

\section{Reconstruction Details}
\label{app:reconstruction}

\subsection{Streams and session state}

All confirmatory Haiku, Opus, and GPT comparisons use canonical seeds
2000--2011. For each seed, a local pseudorandom generator first permutes the
eight procedure families
\texttt{lcg}, \texttt{modpow}, \texttt{continued\_frac},
\texttt{crt\_solve}, \texttt{josephus}, \texttt{quadratic\_map\_mod},
\texttt{xorshift\_steps}, and \texttt{matrix\_power\_mod} across three hot and
five trap roles. Hot classes jointly receive probability 0.85 and traps 0.15,
uniformly within role. The generator samples $T=60$ class identities
i.i.d. and rejection-samples until at least one trap occurs in the first $B$
positions ($B=3$ in the core comparison). Numeric instances are generated at
magnitude 100 for the canonical Haiku conditions. The repaired Opus R3 arm
holds the complete latent stream fixed---family, class ID, role, rate, realized
count, occurrence position, and arrival order---and regenerates only numeric
payloads at magnitude 1000. It does not pin or relocate any family. Structural
projections match the canonical Opus urn artifacts on 12/12 seeds. Class
metadata are hidden; only rendered item or problem text is shown to the model.
Haiku R0, R1, R2, R2c, and R3 reuse the same saved magnitude-100 stream for
each seed.

Each commitment consumes one of $B$ slots and is irreversible. Subsequent
members of a committed class are automatically collected or solved and produce
no new decision. Sessions stop making decisions when the budget is exhausted,
but scoring includes all later members of committed classes. Economic sessions
use the same streams and cross $B\in\{1,3,5\}$ with
$K\in\{0,10,20,24\}$; $K$ changes only the disclosed one-time charge and net
score.

\subsection{Prompt and tool reconstruction}

The artifact retains every rendered user message, assistant response, tool
call, and generated stream. Exact prompt generators and schemas are in
\texttt{urn\_common.py}, \texttt{urn\_session.py},
\texttt{urn\_tool\_session.py}, \texttt{claim\_solver\_session.py},
\texttt{claim\_solver\_code\_session.py}, and
\texttt{session\_state.py}. To remove ambiguity about substitutions, the core
A2 prompts have $T=60$, $N=8$, and the current $B$ inserted into the following
templates.

\paragraph{R0/R1 allocation instructions.}
The system states that 60 items are drawn one at a time; at any draw the model
may KEEP the current item or PASS; at most $B$ items may be kept; keeping an
item immediately collects it and every future item with the same visible
attribute; passing scores zero; already-kept attributes are auto-collected and
do not create another decision; and choices do not change the stream. The A2
disclosure is:
\begin{quote}\small
There are exactly 8 distinct colors in the bag, but you are not told which
colors are common or rare, nor how many balls of each color there are; you
learn the rates only from the draws you have seen.
\end{quote}
R0 ends with the exact response instruction
\begin{quote}\small
For each draw, think briefly (one or two sentences), then end your reply with
EXACTLY one line: \texttt{DECISION: KEEP} or
\texttt{DECISION: PASS}.
\end{quote}
R1 replaces only that instruction with an exactly-one-call requirement. Its
zero-argument functions are \texttt{keep}, described as collecting the current
ball and every future ball of its color while using one keep, and
\texttt{pass}, described as discarding the current ball for zero score. Both
schemas have JSON type \texttt{object} with no properties. The exact R1
response instruction is:
\begin{quote}\small
For each draw, decide by calling EXACTLY ONE tool: keep or pass. Do NOT write a
``DECISION:'' line or any other text describing your choice---the tool call
itself IS your decision. You may think briefly in plain text first if you like,
but you must still call keep or pass.
\end{quote}

\paragraph{R2/R2c allocation instructions.}
Both prompts state that 60 numeric problems arrive one at a time, each has a
hidden type, exactly eight types occur, rates are initially unknown, a solver
may be claimed at most $B$ times, a claim scores the current and all future
problems of that type, skipping scores zero, decisions are final, and the model
must call exactly one of \texttt{claim\_solver} or \texttt{skip\_solver}.
R2 explicitly says not to compute an answer and uses two zero-argument tools.
R2c instead adds the following exact requirement:
\begin{quote}\small
You must provide \texttt{code}: a real, correct Python function
\texttt{def solve(inputs: dict) -> float} that solves problems of this type---it
stands in for every future problem of the same type for the rest of the
session, so it should actually work, not just handle the problem currently
shown.
\end{quote}
The R2c \texttt{claim\_solver} schema therefore has one required string
property, \texttt{code}; \texttt{skip\_solver} remains a zero-argument object.
The prompt does not disclose that correctness is excluded from reward. For an
economic cell, both R0 and R2c add the exact statement that committing costs
$K$ points once, each collected or solved problem remains worth one point, and
declining costs nothing. R3 uses the reusable-script harness's
\texttt{write\_script}, \texttt{run\_script}, and answer tools; the complete
prompt and schema generators are in \texttt{driver.py} and
\texttt{session\_state.py}, and complete resulting transcripts are stored in
\texttt{runs/arm\_a1\_announce\_n-announced/} for Haiku and
\texttt{runs/arm\_a1\_announce\_opus\_n-announced\_canonical-structure/}
for the repaired Opus endpoint.

\subsection{Models and generation settings}

Claude runs use \texttt{claude-haiku-4-5-20251001} and
\texttt{claude-opus-4-8}. GPT runs use
\texttt{gpt-5.4-mini-2026-03-17} and \texttt{gpt-5.6-sol}. Provider-default
temperature is held fixed across paired arms. Claude R1, R2, and R2c use
\texttt{tool\_choice=auto}; forcing a tool call suppressed deliberation in
calibration. GPT uses \texttt{reasoning\_effort=none} in both arms and
\texttt{tool\_choice=auto} for R2c. Maximum output is 512 tokens for R0/R1,
1500 for R2, 2500 for R2c, and 4096 per R3 turn, with a 300,000-token session
cap. Opus R3 sessions stop immediately after the third successful script write;
all 12 terminate on budget exhaustion without hitting the token or per-seed
cost breaker, and contain no malformed or unknown calls. Qwen starts from
\texttt{Qwen/Qwen2.5-Coder-14B-Instruct}; matched final
tool evaluations use q8\_0 base and adapted models with an 8192-token context.

\subsection{Competence gate and mechanical validity}

A GPT model passes the locked abstract competence gate only if all 12 R0
sessions complete, it makes at least 2.5 commitments per seed, no more than
50\% of realized commitments are first-sight, and its mean collection reaches
90\% of the same-information Bayesian comparator. The preregistered
suppression replication additionally requires at least 80\% R2c first-sight
among commitments, a pooled R2c--R0 increase of at least 30 points, and a
positive lower endpoint for the paired bootstrap interval.

Transport or provider exceptions receive at most two additional attempts with
linear 1.5-second backoff. Model-level failures---no decision, multiple
decisive calls, an unknown function, malformed arguments, or an unparsable
response---are not retried; they default to the non-commit action and remain
counted as unresolved. A ladder rung is invalid if unresolved decisions exceed
10\%. Response-model identity, stream hash, malformed and unknown calls,
termination reason, usage, and token-cap hits are checked separately from
behavior. Only a mechanical failure permits a replacement run. For the final
Qwen script evaluation, empty-fence retries are fixed at four in both arms and
all attempts remain logged.

\subsection{Bootstrap estimands}

For each paired comparison, first-sight among commitments is computed within
seed, the second-arm value minus the first-arm value is taken, and paired seeds
are sampled with replacement 10,000 times using bootstrap RNG seed 0. The
reported point estimate is the arithmetic mean of seed differences and the
interval is the percentile interval. A seed with no commitments has an
undefined among-commitments proportion and is omitted only from that
conditional estimand; it remains in commitment, zero-commit, hazard, and
utility summaries. Economic endpoint interactions are computed per seed as
\[
\{h_{R0}(0)-h_{R0}(24)\}
-\{h_{R2c}(0)-h_{R2c}(24)\},
\]
separately for each $B$, and then resampled by seed. No session within a seed
is treated as an independent replicate.

\section{Complete Economic Summaries}
\label{app:economic}

Table~\ref{tab:econ-full} reports the principal behavioral and utility
quantities on the plotted charge surface. Hazard is pooled over eligible
first-sight turns; all other columns are means over the 12 seeds. Zero-commit
incidence is 0\% in every displayed cell.

\begin{table}[H]
\centering
\scriptsize
\caption{Economic surface with timing, commitment, and utility outcomes reported
together. Net is collected or solved items minus $K$ times commitments; regret
is against the exact hindsight net optimum.}
\label{tab:econ-full}
\begin{tabular}{llrrrrr}
\hline
Frame & $(B,K)$ & Hazard (\%) & Commits & Lateness & Net & Regret \\
\hline
R0 & $(1,0)$  & 16.7 & 1.00 & 1.08 & 8.8 & 10.2 \\
R0 & $(1,10)$ & 0.0  & 1.00 & 3.75 & 3.2 & 5.8 \\
R0 & $(1,24)$ & 0.0  & 1.00 & 3.17 & -11.0 & 11.0 \\
R0 & $(3,0)$  & 25.0 & 3.00 & 1.08 & 36.9 & 12.6 \\
R0 & $(3,10)$ & 5.6  & 3.00 & 2.00 & 10.8 & 8.8 \\
R0 & $(3,24)$ & 8.6  & 3.00 & 1.94 & -34.2 & 34.2 \\
R0 & $(5,0)$  & 41.6 & 4.75 & 0.68 & 51.2 & 4.7 \\
R0 & $(5,10)$ & 16.2 & 4.75 & 1.57 & -0.1 & 19.7 \\
R0 & $(5,24)$ & 4.7  & 4.25 & 2.36 & -60.3 & 60.3 \\
\hline
R2c & $(1,0)$  & 100.0 & 1.00 & 0.00 & 9.8 & 9.2 \\
R2c & $(1,10)$ & 100.0 & 1.00 & 0.00 & -0.2 & 9.2 \\
R2c & $(1,24)$ & 100.0 & 1.00 & 0.00 & -14.2 & 14.2 \\
R2c & $(3,0)$  & 94.6 & 3.00 & 0.06 & 32.3 & 17.2 \\
R2c & $(3,10)$ & 100.0 & 3.00 & 0.00 & 2.7 & 16.9 \\
R2c & $(3,24)$ & 97.3 & 3.00 & 0.00 & -38.1 & 38.1 \\
R2c & $(5,0)$  & 100.0 & 5.00 & 0.00 & 50.2 & 5.8 \\
R2c & $(5,10)$ & 100.0 & 5.00 & 0.00 & 0.2 & 19.4 \\
R2c & $(5,24)$ & 98.3 & 5.00 & 0.02 & -69.9 & 69.9 \\
\hline
\end{tabular}
\end{table}

\section{Reward-Learning Details}
\label{app:rl}

The single analyzed run uses QLoRA and per-decision PPO for 20 outer steps,
25 fresh training seeds per step, four rollouts per seed, one epoch per update,
learning rate $6\times10^{-5}$, and rollout temperature 1.2. Training seeds
begin at 9000 and do not overlap evaluation seeds. The critic observes the
actor's ordinary state plus true class-rate features; the actor never receives
those rates. Exact Monte Carlo returns are used without discounting. No
demonstrations enter training.

The unchanged final checkpoint is evaluated on the original ball vocabulary
and on treasure-chest/coin/crest, quiver/arrow/feather, and
cauldron/potion/label reskins. All retain $N=8,T=60,B=3$, the same latent
streams, and free-text keep/pass decisions. A second probe changes only the
decision channel to required zero-argument keep/pass calls. Transfer varies
from 44\% to 97\% first sight across these tool-call reskins; the cauldron
condition accounts for the strongest collapse, with the first three distinct
classes selected by decision turn three in 16 of 24 seeds. We do not assign a
semantic cause because the relevant transcripts contain no explanatory
reasoning. The final reusable-script comparison uses 24 paired seeds
2000--2023, matched q8\_0 quantization, calibrated script success, and four
empty-fence retries in both base and RL-final arms.

\section{Commands and Artifact Manifest}

The core runs are reproduced from the repository root with
\begin{verbatim}
PYTHONPATH=. python -u -m scripts.creator.tool_disposition_benchmark.urn_session \
  --model haiku --announce-n --seeds $(seq 2000 2011)
PYTHONPATH=. python -u -m scripts.creator.tool_disposition_benchmark.urn_tool_session \
  --model haiku --announce-n --seeds $(seq 2000 2011)
PYTHONPATH=. python -u -m scripts.creator.tool_disposition_benchmark.claim_solver_session \
  --model haiku --announce-n --seeds $(seq 2000 2011)
PYTHONPATH=. python -u -m scripts.creator.tool_disposition_benchmark.claim_solver_code_session \
  --model haiku --announce-n --seeds $(seq 2000 2011)
PYTHONPATH=. python -u -m scripts.creator.tool_disposition_benchmark.run_economic_surface
PYTHONPATH=. python -u -m scripts.creator.tool_disposition_benchmark.analyze_economic_surface
\end{verbatim}
Core artifacts are
\texttt{runs/urn\_haiku\_n-announced},
\texttt{runs/urn\_tool\_haiku\_n-announced},
\texttt{runs/claim\_solver\_haiku\_n-announced},
\texttt{runs/claim\_solver\_code\_haiku\_n-announced},
\texttt{runs/arm\_a1\_announce\_n-announced}, and
\texttt{runs/economic\_surface\_haiku}. GPT artifacts are under
\texttt{runs/economic\_surface\_gpt-5.4-mini-2026-03-17} and
\texttt{runs/economic\_surface\_gpt-5.6-sol}; each contains the exact model
identifier, settings, stream, transcript, usage, and per-seed analysis.

\end{document}
